%Two resources useful for abstract writing.
% Guidance of how to write an abstract/summary provided by Nature: https://cbs.umn.edu/sites/cbs.umn.edu/files/public/downloads/Annotated_Nature_abstract.pdf %https://writingcenter.gmu.edu/guides/writing-an-abstract
\chapter*{\center \Large  Abstract}

In recent years Recommender Systems have gained popularity by addressing the problem of consumer choice overload. Conversational Recommender Systems go one step further by incorporating interactivity, enabling active learning and real-time user preference updates to provide more personalised recommendations, closely resembling human conversations. While this research has garnered a lot of attention recently, most existing systems are based on synthetic datasets or predefined templates and rule-based logic. The challenge lies in developing CRS algorithms that can sustain natural dialogues and produce high-quality recommendations.

The primary objective of this thesis is to develop a Conversational Recommendation system that can engage in natural language conversations and provide high-quality recommendations. The system utilises the bot-play approach for training, enabling it to learn from data without relying on synthetic training environments. Specifically, the focus is on optimizing the Information Elicitation module to select questions that yield the highest information gain, leading to high-quality recommendations with minimal conversational steps.

The thesis commences with a comprehensive literature review on Conversational Recommender Systems, delving into the latest research in the field and examining reinforcement learning techniques employed for their training. Subsequently, the recommendation module is discussed in detail, presenting a novel architecture that incorporates an item-to-item approach using a Long Short-Term Memory (LSTM) network with attention mechanism to predict movie ratings. This design allows the module to make precise predictions for new users based on limited preference information acquired through dialogue. To enhance the elicitation of user preferences, the thesis proposes a reinforcement learning technique called bot-play, where a QBot proactively prompts an ABot to provide movie ratings, leveraging the reduction in recommendation model loss as a reward. This enables the model to learn an optimal questioning strategy to maximize the accuracy of user profile representation and recommendation relevance.

The experimental results demonstrate proficiency of the recommendation component in learning item-attribute mappings, enabling the QBot to make accurate rating predictions with only a limited number of answered questions. Furthermore, the trained policy model surpasses the random baseline by effectively capturing user preferences on both synthetic and real-world datasets, showcasing its ability to minimize recommendation model error with minimal conversational turns and reduce loss compared to the baseline. The primary implication of this work lies in the proposed algorithm, which addresses the challenge of developing an efficient active learning strategy for conversational recommendation systems -- a critical aspect often overlooked by existing systems that primarily focus on mimicking human speech. By combining this approach with a natural language interface, it has the potential to yield a practical and valuable conversational system suitable for real-world applications.

%%%%%%%%%%%%%%%%%%%%%%%%%%%%%%%%%%%%%%%%%%%%%%%%%%%%%%%%%%%%%%%%%%%%%%%%%s
~\\[1cm]
\noindent % Provide your key words
\textbf{Keywords:} conversation recommender systems, neural recommender systems, bot-play, information elicitation